problem:  
Let \((M^n,g)\) be a complete noncompact Riemannian manifold with \(\operatorname{Ric}_M \ge 0\), and let \((N^m,h)\) be a compact Riemannian manifold with nonpositive sectional curvature.  
Suppose \(f:(M,g)\to (N,h)\) is a finite-energy harmonic map.  

For \(R>0\), define the rescaled maps \(f_R:(M,g_R=g/R^2)\to(N,h)\) centered at a base point \(p\in M\).  
By the Cheeger–Colding theory, there exists a subsequence \(R_i\to\infty\) such that  
\[
(M,g_{R_i},p)\longrightarrow (C_\infty,o)
\]
in the pointed Gromov–Hausdorff sense, where \(C_\infty\) is a metric cone with vertex \(o\).  
Assume that a subsequence of the energy measures
\[
\mu_R = |df_R|^2\, d\mathrm{vol}_{g_R}
\]
converges weakly to a finite Radon measure \(\mu_\infty\) on \(C_\infty\), and that the maps \(f_{R_i}\) converge weakly in \(W^{1,2}_{\mathrm{loc}}(C_\infty,N)\) to a limit map \(f_\infty\).  

> **Question:**  
> (1) Is the limiting measure \(\mu_\infty\) necessarily *homogeneous* with respect to the cone dilations on \(C_\infty\)?  
>  
> (2) If so, does it follow that the limit map \(f_\infty\) is *homogeneous of degree zero* (i.e. invariant under cone dilations up to isometries of \(N\))?  
>  
> Equivalently, can one expect any finite-energy harmonic map from a Ricci ≥ 0 manifold to become, after blow-down, an asymptotically conical harmonic map on the metric cone at infinity—thereby revealing a new form of *asymptotic homogeneity rigidity*?

Why is it a "good" problem:  
This problem is tightly grounded in current research-level frameworks of geometric analysis and Ricci-limit geometry, yet it is both simple to state and deeply new in focus.  

1. **It reformulates asymptotic rigidity in a natural geometric–analytic language.**  
   Instead of discussing where \(|df|\) attains a maximum, the problem probes the behavior of energy measures under blow-down convergence—an invariant meaningful even on singular Ricci-limit spaces. Conical invariance of the limiting energy measure reflects an analytic analogue of the geometric cone structure of \((M,g)\) at infinity; whether this analytic symmetry must hold is completely nontrivial.

2. **It bridges modern fields.**  
   The question lies at the intersection of  
   - the analysis of harmonic maps and their energy distributions,  
   - the Cheeger–Colding theory for Ricci-nonnegative manifolds, and  
   - the study of weak \(W^{1,2}\)-limits on singular spaces.  
   This is precisely the type of synthesis that drives new discoveries in geometric analysis.  

3. **It targets a fundamental frontier:**  
   A positive result would establish a geometric–analytic *asymptotic rigidity principle*, showing that nonnegative Ricci curvature forces harmonic maps to become conically homogeneous in any blow-down limit.  
   A counterexample would reveal unexpected analytic flexibility of harmonic maps under Ricci nonnegativity, opening new lines of inquiry into the structure of noncollapsed limits and energy concentration phenomena.  

The statement is clean: “Does the blow-down harmonic map inherit the cone’s dilation symmetry?” It is immediately understandable and yet would likely require merging harmonic map regularity, measure convergence, and the delicate geometry of metric cones—exemplifying the “narrow entrance, profound depth” nature of a truly good research problem.